% GENERATED FILE. Edit canonical lessons, then run npm run generate:books.
% canonical-source-hash: fnv1a64:dcda8e5e7aa8d135
% canonical-lessons: BN-C02-naam, BN-C02-amar, BN-C02-amar-naam, BN-C02-alaap, BN-C02-ki, BN-C02-tumi-apni, BN-C02-tomar-naam-ki, BN-C02-practice

\chapter{Introducing Yourself}
\label{ch:introductions}
\hlchaptermodality{2}

\begin{chapteropening}
\textbf{By the end of this chapter:} \emph{I can introduce myself in Bengali, give my name, and understand the other person's name.}

Complete BN-C02-practice, the canonical closing checkpoint, and demonstrate the chapter promise: introduce myself in Bengali, give my name, and understand the other person's name.
\end{chapteropening}

\section[\bn{নাম}]{\bn{নাম} (nām) — \textquotedblleft{}name,\textquotedblright{} a word older than Europe}

\label{lesson:BN-C02-naam}



\begin{quote}
The first word of the introduction — and one of the clearest bridges between India and Europe.
\end{quote}

\subsection*{The letters in this word}
\textbf{\bn{ন}} (\emph{nô}) + \textbf{\bn{া}} (the long-\emph{ā} sign) + \textbf{\bn{ম}} (\emph{mô}) $\to$ \textbf{\bn{নাম}} (\emph{nām}). Here the long-\emph{ā} sign overrides the inherent \emph{ô}, so it is \emph{nām}, not \textquotedblleft{}nôm.\textquotedblright{}

\begin{cousinweb}
\textbf{\bn{নাম}} (\emph{nām}, \textquotedblleft{}name\textquotedblright{}) $\leftarrow$ Sanskrit \textbf{\bn{নামন্}} (\emph{nāman}), from Proto-Indo-European \textbf{*h\textsubscript{3}nómn} (its final \emph{n} syllabic) — the \emph{same} ancient root as English \textbf{name}, Latin \emph{nōmen} (English \textbf{noun}, \textbf{nom}inate), and Greek \emph{onoma}. Bengali \emph{nām} and English \emph{name} are cousins from before either language existed — one of the clearest threads tying the Indian and European branches of one family.
\end{cousinweb}

\subsection*{Your turn}
Say these aloud:

\begin{itemize}
\raggedright
  \item \textquotedblleft{}nām\textquotedblright{}
  \item the family — \emph{nām} / English \emph{name} / Latin \emph{nōmen} / Greek \emph{onoma}
\end{itemize}

\begin{tcolorbox}[breakable,colback=teal!4,colframe=teal!35!black,title={Before you move on}]
What roots is \emph{nām} from, and its English cousins? (Sanskrit \emph{nāman} / PIE \emph{*h\textsubscript{3}nómn}; \textbf{name}, \textbf{noun}, \textbf{nominate}.)
\end{tcolorbox}
\section[\bn{আমার}]{\bn{আমার} (āmār) — \textquotedblleft{}my\textquotedblright{}}

\label{lesson:BN-C02-amar}



\begin{quote}
The little word that makes the name \emph{yours}.
\end{quote}

\subsection*{The letters in this word}
\textbf{\bn{আ}} (the independent long \emph{ā}, used because the word starts with a vowel) + \textbf{\bn{মা}} (\emph{mā}) + \textbf{\bn{র}} (\emph{r}) $\to$ \textbf{\bn{আমার}} (\emph{āmār}).

\begin{cousinweb}
\textbf{\bn{আমার}} (\emph{āmār}, \textquotedblleft{}my\textquotedblright{}) is the possessive of \textbf{\bn{আমি}} (\emph{āmi}, \textquotedblleft{}I\textquotedblright{}), on the first-person stem \textbf{ām-} $\leftarrow$ Sanskrit \emph{asmad} (the \textquotedblleft{}we/I\textquotedblright{} stem). Buried in it is the same ancient first-person \textbf{m-} that runs through English \textbf{me}, \textbf{my}, \textbf{mine} and Latin \emph{mē} — so Bengali's \textquotedblleft{}my\textquotedblright{} and English's \textquotedblleft{}my\textquotedblright{} share, at the root, one first-person sound.
\end{cousinweb}

\begin{grammarlens}[title={Grammar Lens: no gender here}]
Note something Bengali does \emph{not} do: \textbf{\bn{আমার}} never changes for gender. Hindi and Punjabi have \emph{merā} (m.) / \emph{merī} (f.); Bengali has just \textbf{āmār}, for a man or a woman, a masculine or feminine thing. Bengali threw grammatical gender away entirely — you will see this again and again, and it makes Bengali, in this one respect, simpler than its western cousins.
\end{grammarlens}

\subsection*{Your turn}
Say these aloud:

\begin{itemize}
\raggedright
  \item \textquotedblleft{}āmār\textquotedblright{}
  \item its owner (\emph{āmi}, \textquotedblleft{}I\textquotedblright{})
  \item does it change for a man vs. a woman? (No — Bengali has no gender)
\end{itemize}

\begin{tcolorbox}[breakable,colback=teal!4,colframe=teal!35!black,title={Before you move on}]
What is \emph{āmār} the possessive of, and how does it differ from Hindi's \emph{merā}? (\emph{Āmi}, \textquotedblleft{}I\textquotedblright{}; \emph{āmār} never changes for gender, where \emph{merā/merī} does.)
\end{tcolorbox}
\section[\bn{আমার} \bn{নাম} …]{\bn{আমার} \bn{নাম} … (āmār nām …) — \textquotedblleft{}my name is…,\textquotedblright{} with no \textquotedblleft{}is\textquotedblright{}}

\label{lesson:BN-C02-amar-naam}



\begin{quote}
Assemble the two atoms — and notice a whole word English needs that Bengali drops.
\end{quote}

\subsection*{The two words — assembled}
\textbf{\bn{আমার}} (my) + \textbf{\bn{নাম}} (name) + your name $\to$ \textbf{\bn{আমার} \bn{নাম} …} = \textquotedblleft{}my name …\textquotedblright{} $\to$ \textquotedblleft{}\textbf{my name is…}\textquotedblright{} \emph{Āmār nām Arun.} $\to$ \textquotedblleft{}My name is Arun.\textquotedblright{}

\begin{grammarlens}[title={Grammar Lens: Bengali drops \textquotedblleft{}is\textquotedblright{} in the present}]
Look at what is missing: there is \textbf{no verb}. Bengali sets \textquotedblleft{}my name\textquotedblright{} beside \textquotedblleft{}Arun\textquotedblright{} — \emph{āmār nām Arun} — and the \textquotedblleft{}is\textquotedblright{} is simply understood. This is the \textbf{zero copula}: for present-tense equational sentences (X is Y), Bengali needs no \textquotedblleft{}to be\textquotedblright{} at all. (It is the same habit Tamil has — and unlike Hindi's \emph{merā nām … hai}, which keeps the \emph{hai}.) Bengali only brings back a verb when it talks about a \emph{state} or \emph{place} — you'll see that in Chapter 3.
\end{grammarlens}

\subsection*{Your turn}
Say these aloud:

\begin{itemize}
\raggedright
  \item \textquotedblleft{}āmār nām …\textquotedblright{} with your name
  \item notice — no word for \textquotedblleft{}is\textquotedblright{}
\end{itemize}

\begin{tcolorbox}[breakable,colback=teal!4,colframe=teal!35!black,title={Before you move on}]
Give your name in Bengali, and say what it leaves out that English requires. (\emph{Āmār nām ….}; the verb \textquotedblleft{}is\textquotedblright{} — the zero copula.)
\end{tcolorbox}
\section[\bn{আলাপ} \bn{করে} \bn{ভালো} \bn{লাগলো}]{\bn{আলাপ} \bn{করে} \bn{ভালো} \bn{লাগলো} (ālāp kore bhālo lāglo) — \textquotedblleft{}pleased to meet you\textquotedblright{}}

\label{lesson:BN-C02-alaap}



\begin{quote}
The warm close of an introduction.
\end{quote}

\subsection*{The letters in this word}
The key word is \textbf{\bn{আলাপ}} (\emph{ālāp}): \textbf{\bn{আ}} + \textbf{\bn{লা}} (\emph{lā}) + \textbf{\bn{প}} (\emph{p}).

\begin{cousinweb}
\textbf{\bn{আলাপ} \bn{করে} \bn{ভালো} \bn{লাগলো}} = \textbf{\bn{আলাপ}} (\emph{ālāp}, \textquotedblleft{}acquaintance, conversation\textquotedblright{}) + \textbf{\bn{করে}} (\emph{kore}, \textquotedblleft{}having done\textquotedblright{}) + \textbf{\bn{ভালো} \bn{লাগলো}} (\emph{bhālo lāglo}, \textquotedblleft{}{[}it{]} felt good\textquotedblright{}) — literally \textquotedblleft{}\textbf{having made your acquaintance, {[}it{]} felt good}.\textquotedblright{} The key word \textbf{\bn{আলাপ}} (\emph{ālāp}) is \textbf{Sanskrit} — \emph{ālāpa} (\textquotedblleft{}conversing, addressing,\textquotedblright{} from \emph{ā-} \textquotedblleft{}toward\textquotedblright{} + \emph{lap} \textquotedblleft{}to speak\textquotedblright{}), the same \emph{ālāp} that names the slow opening section of a rāga in Indian classical music, where the musician \textquotedblleft{}converses\textquotedblright{} with the notes. So even a hello leans on Sanskrit's word for the meeting of two voices.
\end{cousinweb}

\subsection*{Your turn}
Say these aloud:

\begin{itemize}
\raggedright
  \item \textquotedblleft{}ālāp kore bhālo lāglo\textquotedblright{}
  \item what \emph{ālāp} means, and where else you've heard it (acquaintance/conversation; the \emph{ālāp} of a rāga)
\end{itemize}

\begin{tcolorbox}[breakable,colback=teal!4,colframe=teal!35!black,title={Before you move on}]
What does the phrase literally say, and what is \emph{ālāp}? (\textquotedblleft{}Having made {[}your{]} acquaintance, it felt good\textquotedblright{}; Sanskrit \emph{ālāpa}, \textquotedblleft{}conversing\textquotedblright{} — as in a rāga's opening.)
\end{tcolorbox}
\section[\bn{কি}]{\bn{কি} (ki) — \textquotedblleft{}what\textquotedblright{}}

\label{lesson:BN-C02-ki}



\begin{quote}
The question-word — and it starts with the \textbf{k-} that heads nearly every Bengali question.
\end{quote}

\subsection*{The letters in this word}
\textbf{\bn{ক}} (\emph{kô}) + \textbf{\bn{ি}} (the \emph{i}-sign, written \textbf{before} the consonant but read after) $\to$ \textbf{\bn{কি}} (\emph{ki}).

\begin{cousinweb}
\textbf{\bn{কি}} (\emph{ki}, \textquotedblleft{}what\textquotedblright{}) $\leftarrow$ Sanskrit \emph{kim} / the interrogative stem \textbf{ka-}, from PIE \textbf{*k\textsuperscript{w}o-} — the same question-root behind English \textbf{wh-} words (\textbf{what}, \textbf{who}, \textbf{which}) and Latin \emph{qu-}. Bengali's questions nearly all begin with this \textbf{k-}: \emph{ki} (what), \emph{ke} (who), \emph{kôbe} (when), \emph{kōthāy} (where), \emph{kēmon} (how). Learn the \emph{k-} and you can hear a question coming.
\end{cousinweb}

\subsection*{Your turn}
Say these aloud:

\begin{itemize}
\raggedright
  \item \textquotedblleft{}ki\textquotedblright{}
  \item the k- question family — \emph{ki, ke, kôbe, kōthāy, kēmon}
\end{itemize}

\begin{tcolorbox}[breakable,colback=teal!4,colframe=teal!35!black,title={Before you move on}]
What root do Bengali's question-words share, and its English family? (The interrogative \emph{k-}, PIE \emph{*k\textsuperscript{w}o-}; the \emph{wh-} words.)
\end{tcolorbox}
\section[\bn{তুমি} / \bn{আপনি}]{\bn{তুমি} / \bn{আপনি} (tumi / āpni) — \textquotedblleft{}you,\textquotedblright{} three ways}

\label{lesson:BN-C02-tumi-apni}



\begin{quote}
To ask a name you need \textquotedblleft{}you\textquotedblright{} — and Bengali grades it \textbf{three} ways.
\end{quote}

\subsection*{The letters in this word}
\textbf{\bn{তুমি}} (\emph{tumi}): \textbf{\bn{তু}} (\emph{tu}) + \textbf{\bn{মি}} (\emph{mi}). \textbf{\bn{আপনি}} (\emph{āpni}): \textbf{\bn{আ}} + \textbf{\bn{প}} + \textbf{\bn{নি}}.

\begin{cousinweb}
\textbf{\bn{তুমি}} (\emph{tumi}, familiar \textquotedblleft{}you\textquotedblright{}) $\leftarrow$ Sanskrit \emph{tvam}, from PIE \textbf{*tū} — the same root as English archaic \textbf{thou}, Latin \emph{tū}, French \emph{tu}. Bengali grades respect in \textbf{three} levels:

\begin{itemize}
\raggedright
  \item \textbf{\bn{তুই}} (\emph{tui}) — intimate: close friends, children, family.
  \item \textbf{\bn{তুমি}} (\emph{tumi}) — familiar: friends, juniors, warmth.
  \item \textbf{\bn{আপনি}} (\emph{āpni}) — respectful: elders, strangers, formality.
\end{itemize}

(This is like Hindi's \emph{tū / tum / āp}, and finer than the two-way split of the Romance tracks.) For a first meeting, use \textbf{\bn{আপনি}}.
\end{cousinweb}

\subsection*{Your turn}
Say these aloud:

\begin{itemize}
\raggedright
  \item the three — \emph{tui} (intimate) / \emph{tumi} (familiar) / \emph{āpni} (respectful)
  \item the English cousin of \emph{tumi} (\textbf{thou})
  \item which you use on first meeting (\emph{āpni})
\end{itemize}

\begin{tcolorbox}[breakable,colback=teal!4,colframe=teal!35!black,title={Before you move on}]
Name Bengali's three \textquotedblleft{}you\textquotedblright{}s, and which English pronoun \emph{tumi} is cousin to. (\emph{Tui / tumi / āpni}; archaic \textbf{thou}.)
\end{tcolorbox}
\section[\bn{তোমার} \bn{নাম} \bn{কি}?]{\bn{তোমার} \bn{নাম} \bn{কি}? (tomār nām ki?) — \textquotedblleft{}what's your name?\textquotedblright{}}

\label{lesson:BN-C02-tomar-naam-ki}



\begin{quote}
Turn the introduction around — ask, instead of tell.
\end{quote}

\subsection*{The two words — assembled}
\textbf{\bn{তোমার}} (\emph{tomār}, \textquotedblleft{}your,\textquotedblright{} from \emph{tumi}) + \textbf{\bn{নাম}} (name) + \textbf{\bn{কি}} (what) $\to$ literally \textquotedblleft{}\textbf{your name what?}\textquotedblright{} And note — still \textbf{no \textquotedblleft{}is.\textquotedblright{}} The zero copula holds here too: Bengali simply lays \textquotedblleft{}your name\textquotedblright{} beside \textquotedblleft{}what,\textquotedblright{} and the question is complete.

\begin{grammarlens}[title={Grammar Lens: answering, and the respectful form}]
Answer with the sentence you built: \emph{Āmār nām Arun.} The respectful version swaps \emph{tomār} (from \emph{tumi}) for \textbf{\bn{আপনার}} (\emph{āpnār}, \textquotedblleft{}your,\textquotedblright{} from \emph{āpni}): \emph{āpnār nām ki?} — the one to use with a stranger.
\end{grammarlens}

\subsection*{Your turn}
Say these aloud:

\begin{itemize}
\raggedright
  \item the familiar question — \emph{tomār nām ki?}
  \item the respectful version — \emph{āpnār nām ki?}
  \item answer it — \emph{āmār nām …}
\end{itemize}

\begin{tcolorbox}[breakable,colback=teal!4,colframe=teal!35!black,title={Before you move on}]
Ask \textquotedblleft{}what's your name?\textquotedblright{} respectfully, and note what's still missing. (\emph{Āpnār nām ki?}; there is no word for \textquotedblleft{}is\textquotedblright{} — the zero copula.)
\end{tcolorbox}
\section[Practice]{Chapter 2 — The introduction exchange}

\label{lesson:BN-C02-practice}



\begin{quote}
No new words. The whole introduction, assembled from atoms you now own.
\end{quote}

\subsection*{The exchange}
Say these aloud:

\begin{itemize}
\raggedright
  \item \textquotedblleft{}My name is Mira.\textquotedblright{} — \emph{āmār nām Mira.}
  \item \textquotedblleft{}What's your name?\textquotedblright{} (respectful) — \emph{āpnār nām ki?}
  \item \textquotedblleft{}My name is Arun.\textquotedblright{} — \emph{āmār nām Arun.}
  \item \textquotedblleft{}Pleased to meet you.\textquotedblright{} — \emph{ālāp kore bhālo lāglo.}
\end{itemize}

\begin{cousinweb}
Say these aloud:

\begin{itemize}
\raggedright
  \item \textquotedblleft{}name,\textquotedblright{} \textquotedblleft{}my\textquotedblright{} (\emph{nām}, \emph{āmār})
  \item the three \textquotedblleft{}you\textquotedblright{} levels (\emph{tui / tumi / āpni})
  \item \textquotedblleft{}what,\textquotedblright{} and the word Bengali leaves out (\emph{ki}; there is no \textquotedblleft{}is\textquotedblright{})
\end{itemize}
\end{cousinweb}

\begin{grammarlens}[title={What you've built — roots you now carry}]
\begin{itemize}
\raggedright
  \item \emph{nām} $\leftarrow$ Sanskrit \emph{nāman} $\to$ English \textbf{name}
  \item \emph{āmār} $\leftarrow$ the first-person \emph{ām-}; no gender (unlike Hindi \emph{merā/merī})
  \item \emph{ki} $\leftarrow$ \emph{ka-} $\to$ English \textbf{what}; \emph{tumi} $\leftarrow$ \emph{*tū} $\to$ \textbf{thou}
  \item \emph{ālāp} $\leftarrow$ Sanskrit \emph{ālāpa}, \textquotedblleft{}conversing\textquotedblright{}
\end{itemize}
\end{grammarlens}

\begin{tcolorbox}[breakable,colback=teal!4,colframe=teal!35!black,title={Before you move on}]
Introduce yourself and ask a stranger's name in Bengali — and note the one word you never say. (\emph{Āmār nām …. Āpnār nām ki?} — no word for \textquotedblleft{}is.\textquotedblright{})
\end{tcolorbox}
